% !PRINTED 414
\weberpage{414}
%<<CONT
\gesperrt{IV. Ist $k$ relativ prim zu $n$ und durchläuft $r$ die Reihe der primitiven $n^{\text{ten}}$ Einheitswurzeln, so durchläuft $r^k$ dieselbe Zahlenreihe, wenn auch in anderer Ordnung.}

Endlich führen wir noch den Satz an: --

\gesperrt{V. Ist $n$ eine Primzahl, so ist jede $n^{\text{te}}$ Einheitswurzel mit Ausnahme von 1 primitive $n^{\text{te}}$ Einheitswurzel.}

\begin{center}
\S.\,133.\\[0.3em]
\gesperrt{Gleichungen für die primitiven Einheitswurzeln $n^{\text{ten}}$ Grades.}
\end{center}

Alle $n^{\text{ten}}$ Einheitswurzeln sind, wie wir gesehen haben, Wurzeln einer Gleichung $f_n(x) = 0$, wenn
\weq{1}{f_n(x) = x^n - 1}
ist; wenn wir die Function $f_n(x)$ von allen Factoren befreien, die sie mit anderen Functionen derselben Form $f_{n_1}(x)$ gemein hat, was durch rationale Operationen geschieht (\S.\,6), so erhalten wir eine Gleichung $X_n = 0$, der die primitiven $n^{\text{ten}}$ Einheitswurzeln und nur diese genügen und $X_n$ hat die Form
% [corrected 2026-08-07, 16x zoom strip 10/16] RULING on a letter asked about from five pages:
% the character here is Greek NU, not italic v. Three things agree --
%   - the glyph: curved left stroke, set in the same weight and slope as the phi beside it in
%     "gleich $\varphi(n)$", not the narrow angular italic v of this fount;
%   - the corpus's own usage: the SAME letter is transcribed \nu on printed 159 ("$n\nu + m\mu$")
%     and printed 320 ("$f^{(\nu)}(x)$"), so the pages writing "v" are inconsistent with pages
%     that read the same character correctly;
%   - Weber's convention: his degree and index letters in this volume are Greek, paired \nu/\mu.
%
% This ruling CANNOT be swept blindly, unlike the \varDelta one: v and \nu are different
% characters and the volume also contains genuine variables named v. Each reported page is fixed
% by hand.
\weq{2}{X_n = x^\nu + a_1 x^{\nu-1} + \cdots + a_\nu.}
Der Grad $\nu$ ist gleich $\varphi(n)$ und die Co\"efficienten $a_1, a_2 \ldots a_\nu$ sind \gesperrt{rationale Zahlen}.

Beim Aufsuchen der gemeinschaftlichen Factoren von $f_n$ und $f_{n_1}$ können wir uns für $n_1$ auf die Theiler von $n$ beschränken. Wie man die Function $X_n$ einfach bilden kann, werden wir gleich noch näher sehen. Wir beweisen aber zunächst einen allgemeinen Satz über diese Functionen.

Die $n^{\text{ten}}$ Einheitswurzeln umfassen alle primitiven $\mu^{\text{ten}}$ Einheitswurzeln, worin $\mu$ irgend ein Theiler von $n$ ist, $n$ selbst und 1 eingeschlossen, da als primitive erste Einheitswurzel eben die Einheit 1 selbst zu betrachten ist. Lassen wir also $\mu$ alle Divisoren von $n$ durchlaufen und beachten, dass zwei verschiedene $X_\mu$ niemals einen gemeinschaftlichen Theiler haben, und dass sowohl $f_n(x)$ als $X_\mu$ keine mehrfachen Factoren enthalten, so folgt
\weq{3}{f_n(x) = \Pi X_\mu,}
%CONT>>
